%\begin{center}
%\large \bf \runtitle
%\end{center}
%\vspace{1cm}
\chapter*{\runtitle}

\noindent Antimicrobial resistance in \textit{Acinetobacter baumannii} is a critical clinical challenge because of the frequency of healthcare-associated infections and the limited availability of effective therapeutic options. Identifying novel antibacterial targets requires the integration of evidence reflecting their functional relevance, conservation, and potential for pharmacological intervention. This thesis develops a multicriteria computational workflow to prioritize potential therapeutic targets in clinical strains of \textit{A.~baumannii}. The approach integrates pangenome analysis to identify conserved genes, published experimental evidence of genetic essentiality, functional annotation, involvement in metabolic pathways, and assessment of homology to human proteins. Top-ranked candidates are further characterized through the analysis of protein structures and binding pockets to estimate their ligandability. Integrating these evidence layers produces a reproducible, interpretable, and traceable ranking that makes the strengths and limitations of each candidate explicit. This work provides a strategy to narrow the target search space and guide subsequent experimental validation and rational development of new antibacterials against \textit{A.~baumannii}.

\bigskip

\noindent\textbf{Keywords:} \textit{Acinetobacter baumannii}, antimicrobial resistance, therapeutic targets, genetic essentiality, pangenome, ligandability.
